\documentclass{otscience}

% Also verifies that the workbook compatibility layer can be loaded
% without redefining interfaces already supplied by otscience.cls.
% ============================================================
% Shared setup for OT Science modular workbook parts
% Place these files in the same folder as your fixed otscience.cls
% ============================================================

\makeatletter
\makeatother
\usepackage{multicol}
\usepackage{longtable}
\usepackage{makecell}
\usepackage{pdflscape}
\usepackage{bm}
\providecommand{\blank}[1][3cm]{\rule{#1}{0.4pt}}
\providecommand{\bigblank}{\rule{7cm}{0.4pt}}
\providecommand{\componentblank}{\blank[2.5cm]}
\providecommand{\R}{\mathbb{R}}
\providecommand{\C}{\mathbb{C}}
\providecommand{\ii}{\mathrm{i}}
\providecommand{\ee}{\mathrm{e}}
\providecommand{\dd}{\mathrm{d}}
\providecommand{\grad}{\nabla}
\providecommand{\divergence}{\nabla\!\cdot}
\providecommand{\curl}{\nabla\!\times}
\providecommand{\lap}{\nabla^2}
\providecommand{\ihat}{\hat{\mathbf{i}}}
\providecommand{\jhat}{\hat{\mathbf{j}}}
\providecommand{\khat}{\hat{\mathbf{k}}}
\providecommand{\er}{\hat{\mathbf{e}}_r}
\providecommand{\etheta}{\hat{\mathbf{e}}_\theta}
\providecommand{\ephi}{\hat{\mathbf{e}}_\phi}
\providecommand{\erho}{\hat{\boldsymbol{\rho}}}
\providecommand{\vA}{\mathbf{A}}
\providecommand{\vB}{\mathbf{B}}
\providecommand{\va}{\mathbf{a}}
\providecommand{\vb}{\mathbf{b}}
\providecommand{\vc}{\mathbf{c}}
\providecommand{\vx}{\mathbf{x}}
\providecommand{\vr}{\mathbf{r}}
\providecommand{\matA}{\mathbf{A}}
\providecommand{\matB}{\mathbf{B}}
\providecommand{\tr}{\operatorname{tr}}
\providecommand{\cof}{\operatorname{cof}}
\providecommand{\dev}{\operatorname{dev}}
\providecommand{\diag}{\operatorname{diag}}
\providecommand{\questionline}[1][5cm]{\rule{#1}{0.4pt}}
\setlength{\parindent}{0pt}
\setlength{\parskip}{4pt}
\renewcommand{\arraystretch}{1.25}

% Fallbacks in case the current otscience.cls has not defined nosplit boxes.
\makeatletter
\@ifundefined{formulaboxnosplit}{\newenvironment{formulaboxnosplit}[1][Formula]{\begin{formulabox}[#1]}{\end{formulabox}}}{}
\@ifundefined{definitionboxnosplit}{\newenvironment{definitionboxnosplit}[1][Definition]{\begin{definitionbox}[#1]}{\end{definitionbox}}}{}
\@ifundefined{summaryboxnosplit}{\newenvironment{summaryboxnosplit}[1][Summary]{\begin{summarybox}[#1]}{\end{summarybox}}}{}
\@ifundefined{warningboxnosplit}{\newenvironment{warningboxnosplit}[1][Warning]{\begin{warningbox}[#1]}{\end{warningbox}}}{}
\@ifundefined{exampleboxnosplit}{\newenvironment{exampleboxnosplit}[1][Example]{\begin{examplebox}[#1]}{\end{examplebox}}}{}
\makeatother

\newenvironment{practicebox}[1][Quick Drills and Practice]{\begin{examplebox}[#1]}{\end{examplebox}}
\newenvironment{practiceboxnosplit}[1][Quick Drills and Practice]{\begin{exampleboxnosplit}[#1]}{\end{exampleboxnosplit}}

\newcommand{\standalonetitle}[3]{%
    \title{\textbf{#1}\\\large #2}
    \author{OT Science Notes}
    \date{#3}
}


\standalonetitle
    {OT Science Semantic-Box Compatibility Test}
    {Phase 2 Regression Coverage}
    {4 August 2026}

\begin{document}

\maketitle

\section{Breakable semantic boxes}

\begin{definitionbox}
    The default title should be \textbf{Definition}, with the established
    blue frame.
\end{definitionbox}

\begin{theorembox}
    The default title should be \textbf{Theorem}, with the established
    purple frame. This remains an unnumbered visual box.
\end{theorembox}

\begin{lawbox}
    The default title should be \textbf{Law / Principle}, with the
    established green frame.
\end{lawbox}

\begin{formulabox}
    The default title should be \textbf{Formula}, with the established
    teal frame.
    \[
        E=mc^2.
    \]
\end{formulabox}

\begin{derivationbox}
    The default title should be \textbf{Derivation}, with the established
    orange frame.
\end{derivationbox}

\begin{examplebox}
    The default title should be \textbf{Worked Example}, with the
    established green frame.
\end{examplebox}

\begin{warningbox}
    The default title should be \textbf{Common Mistake}, with the
    established red frame.
\end{warningbox}

\begin{intuitionbox}
    The default title should be \textbf{Intuition}, with the established
    yellow frame.
\end{intuitionbox}

\begin{summarybox}
    The default title should be \textbf{Key Summary}, with the established
    dark frame.
\end{summarybox}

\begin{experimentbox}
    The default title should be \textbf{Experiment / Observation}, with
    the established purple frame.
\end{experimentbox}

\section{Non-splitting semantic boxes}

\begin{definitionboxnosplit}
    This definition box should remain intact.
\end{definitionboxnosplit}

\begin{theoremboxnosplit}
    This theorem box should remain intact and unnumbered.
\end{theoremboxnosplit}

\begin{lawboxnosplit}
    This law box should remain intact.
\end{lawboxnosplit}

\begin{formulaboxnosplit}
    This formula box should remain intact.
    \[
        a^2+b^2=c^2.
    \]
\end{formulaboxnosplit}

\begin{derivationboxnosplit}
    This derivation box should remain intact.
\end{derivationboxnosplit}

\begin{exampleboxnosplit}
    This worked-example box should remain intact.
\end{exampleboxnosplit}

\begin{warningboxnosplit}
    This warning box should remain intact.
\end{warningboxnosplit}

\begin{intuitionboxnosplit}
    This intuition box should remain intact.
\end{intuitionboxnosplit}

\begin{summaryboxnosplit}
    This summary box should remain intact.
\end{summaryboxnosplit}

\begin{experimentboxnosplit}
    This experiment box should remain intact.
\end{experimentboxnosplit}

\section{Custom titles}

\begin{formulabox}[Custom Formula Title]
    The optional argument should replace the default title.
    \[
        \nabla\cdot\mathbf{A}
        =
        \frac{\partial A_x}{\partial x}
        +
        \frac{\partial A_y}{\partial y}
        +
        \frac{\partial A_z}{\partial z}.
    \]
\end{formulabox}

\begin{warningboxnosplit}[Custom Warning Title]
    The optional argument should replace the default title while the box
    retains its established colour and non-splitting behaviour.
\end{warningboxnosplit}

\section{Workbook compatibility wrappers}

\begin{practicebox}
    This box is created through the workbook's breakable
    \texttt{practicebox} wrapper. It should retain the visual appearance
    of an example box.
\end{practicebox}

\begin{practiceboxnosplit}
    This box is created through the workbook's non-splitting
    \texttt{practiceboxnosplit} wrapper. It should remain intact.
\end{practiceboxnosplit}

\clearpage
\section{Breakable page-boundary test}

The following space deliberately positions a long formula box near a
page boundary.

\rule{0pt}{0.48\textheight}

\begin{formulabox}[Breakable Formula Box]
    \foreach \n in {1,...,14}{%
        Test paragraph \n. This repeated material verifies that the
        breakable semantic box can continue onto another page without
        clipping, overlapping the footer, or losing its frame.\par
    }
\end{formulabox}

\clearpage
\section{Non-splitting page-boundary test}

The following space deliberately leaves insufficient room for the
complete box.

\rule{0pt}{0.72\textheight}

\begin{formulaboxnosplit}[Non-Splitting Formula Box]
    This material should move together to the following page when there
    is insufficient space. The title, contents, background, and frame
    must remain intact.
    \[
        \int_a^b f(x)\,dx
        \approx
        \frac{b-a}{2}\bigl(f(a)+f(b)\bigr).
    \]
\end{formulaboxnosplit}

\end{document}